\section{BH4: Order behaviour and participation}\label{sec:bh4}\label{sec:activity}

The fourth hypothesis concerns how participants trade in the window: whether they cancel more,
conceal more, trade with more urgency, and whether the mix of participants changes. Its primary
measure is the share of orders cancelled within one second of entry, the fleeting orders that a
strategy of placing and withdrawing quotes would produce.

\subsection{Cancellation}

The share of orders cancelled within one second rises from \LifePhantomRateControl{}\% to
\LifePhantomRateExpiry{}\% ($d = \HTwoEightEffectSizeCohenD{}$, $p = \HTwoEightPValue{}$), and
the primary test of BH4 is rejected. The rise is specific to securities with derivatives:
relative to the placebo group it is \DidPhantomorderratePlacebo{} larger
($p = \DidPhantomorderratePlaceboP{}$). Cancellation counts rise sharply, which is the largest
effect among all the measures ($d = \HEightEffectSizeCohenD{}$), but cancellations per entry do
not ($p = \HSevenPValue{}$): entries rise as fast as cancellations. The window attracts more
order submission of every kind, and more of the orders it attracts are fleeting, rather than a
disproportionate amount of placing and withdrawing.

\subsection{Concealment}

The share of resting size that is not displayed is high on every session, at
\ConcealedDepthShareControl{}\% on control sessions, and does not change on expiry. Concealed
orders become somewhat larger relative to displayed ones, from \ConcealedSizeMultipleControl{}
to \ConcealedSizeMultipleExpiry{} times their size, and are replenished slightly less often per share
traded. Fewer, larger concealed orders worked more slowly through the window describe the
execution of large expiring positions rather than an attempt to hide an accumulation aimed at
the settlement price.

More entered orders carry a disclosed quantity below their total
($p = \HNinePValue{}$), while the concealed share of submitted volume does not change
($p = \CFourPValue{}$).

\subsection{Urgency}

Immediate-or-cancel orders are used no more often by algorithmic traders on expiry days
($p = \HSixPValue{}$), and their use does not accelerate in the final five minutes
($p = \HOneOnePValue{}$). Trading in the window is not more urgent.

\subsection{Participation}

Proprietary volume in the window rises on expiry sessions ($p = \HThreePValue{}$), while
algorithmic volume does not change significantly ($p = \HFivePValue{}$). Across the year the
shift in participation is visible in the composition of window volume.

Figure~\ref{fig:dailyshare} shows the same quantity session by session across the year. The
largest spikes fall on month ends and index rebalancing dates, and are shared by the placebo
group; the difference between the two on monthly expiries is smaller than these spikes, and
is identified by the regression rather than visible session by session.

\begin{figure}[htbp]
  \centering
  \includegraphics[width=\textwidth]{daily_share.pdf}
  \caption{Share of each session's volume traded in the settlement window, averaged over the
    securities with derivatives and over the placebo group, for every session of the year.
    Vertical lines mark the monthly expiries.}
  \label{fig:dailyshare}
\end{figure}

Institutional participation in the window rises on monthly expiries. In the
liquid group, custodian participants are on at least one side of
\TapeMeanLiquidNormalCustodian{}\% of window volume on ordinary sessions and of
\TapeMeanLiquidMonthlyexpiryCustodian{}\% on monthly expiries, and on
\TapeMeanLiquidMonthendCustodian{}\% on month ends; the participation of other clients falls
correspondingly. Figure~\ref{fig:participants} shows the composition for each group. This is
the pattern expected from institutions closing, rolling or benchmarking positions to the
closing price, and it is strongest where institutional ownership is concentrated.

\begin{figure}[htbp]
  \centering
  \includegraphics[width=\textwidth]{participants.pdf}
  \caption{Share of window volume in which each participant category is on at least one side
    of the trade, by group and type of session. The shares do not sum to one hundred, since a
    trade between two categories counts for both.}
  \label{fig:participants}
\end{figure}

